\documentclass[10pt]{article}

\usepackage[letterpaper,margin=0.58in]{geometry}
\usepackage{amsmath,amssymb}
\usepackage{enumitem}
\usepackage{xcolor}
\usepackage{fancyhdr}

\setlength{\parindent}{0pt}
\setlength{\parskip}{0pt}

\pagestyle{fancy}
\fancyhf{}
\renewcommand{\headrulewidth}{0pt}
\cfoot{\thepage}

\definecolor{sectiongray}{gray}{0.935}

\newcommand{\sectionbar}[1]{%
  \vspace{0.55em}%
  \noindent\colorbox{sectiongray}{%
    \parbox{\dimexpr\textwidth-2\fboxsep\relax}{%
      \textbf{\large\strut #1\strut}%
    }%
  }%
  \par\vspace{0.47em}%
}

\newcommand{\answerbar}[1]{%
  \vspace{0.55em}%
  \noindent\colorbox{sectiongray}{%
    \parbox{\dimexpr\textwidth-2\fboxsep\relax}{%
      \strut\textbf{#1}\strut
    }%
  }%
  \par\vspace{0.50em}%
}

\newlist{parts}{enumerate}{1}
\setlist[parts]{
  label=(\alph*),
  leftmargin=2.1em,
  itemsep=0.45em,
  topsep=0.25em,
  parsep=0pt
}

\begin{document}

% ================================================================
% PAGE 1
% ================================================================

\begin{center}
    {\Large\bfseries Math 2650 -- Quiz 03 Concept Practice}\\[0.28em]
    \small Choose the section matching the skill you want to practice.
    Answers are on the back.
\end{center}

\vspace{0.07em}

% ================================================================
% A
% ================================================================

\sectionbar{A. Motion Along a Space Curve}

An object moves along the curve
\[
\mathbf r(t)
=
\left\langle
t^3+2t,\;
4\sqrt{t+1},\;
2\ln(t+1)
\right\rangle.
\]

\begin{parts}

\item At $t=0$ the object's engine shuts off, after which it travels
with constant velocity $\mathbf r'(0)$. Find its position two units
of time after the engine shuts off.

\item Find the angle between the object's velocity and acceleration
vectors at $t=0$.

\item Find an equation of the plane through $\mathbf r(0)$ that is
perpendicular to the curve at $t=0$.

\end{parts}

\vspace{0.12em}

% ================================================================
% B
% ================================================================

\sectionbar{B. Tangent Lines and Intersecting Curves}

Consider the curve
\[
\mathbf r(t)
=
\langle
2t+1,\;
t^2-3,\;
-t+4
\rangle.
\]

\begin{parts}

\item Find all values of $t$ for which the tangent line to
$\mathbf r(t)$ passes through the point
\[
Q=(5,0,2).
\]

\item A second curve is
\[
\mathbf s(u)
=
\langle
3+u,\;
-2+2u+u^2,\;
3+2u-u^2
\rangle.
\]
The curves intersect at $\mathbf r(1)=\mathbf s(0)$. Let $\theta$
be the angle between the tangent vectors to the two curves at their
point of intersection. Find $\cos\theta$.

\end{parts}

\vspace{0.12em}

% ================================================================
% C
% ================================================================

\sectionbar{C. Quadric Surfaces and Cross Sections}

Consider the surface
\[
4x^2+9y^2-8x+36y-12z+4=0.
\]

\begin{parts}

\item Rewrite the equation in standard form. Identify the quadric
surface, state its axis of symmetry, and give the direction in which
it opens.

\item Identify the conic section obtained by intersecting the surface
with each of the planes
\[
x=1,\qquad y=-2,\qquad z=0.
\]

\item Find the area of the cross section obtained when $z=0$.

\end{parts}

\vspace{0.12em}

% ================================================================
% D
% ================================================================

\sectionbar{D. Spherical and Cartesian Descriptions of a Surface}

A surface is given in spherical coordinates by
\[
\rho
=
\frac{8\sin\phi\sin\theta}
{1+3\cos^2\phi}.
\]

\begin{parts}

\item Rewrite the equation in Cartesian coordinates.

\item Put the Cartesian equation into standard form and identify the
quadric surface.

\item Give its center and semiaxis lengths.

\end{parts}

% ================================================================
% PAGE 2: ANSWER KEY
% ================================================================

\newpage

\begin{center}
    {\Large\bfseries Answer Key}\\[0.28em]
    \small Final answers only
\end{center}

\vspace{0.30em}

% ================================================================
% A ANSWERS
% ================================================================

\answerbar{A. Motion Along a Space Curve}

\begin{parts}

\item
\[
\boxed{(4,8,4)}
\]

\item
\[
\boxed{\theta=\arccos\left(-\frac{\sqrt{15}}5\right)}
\]

\item
\[
\boxed{x+y+z=4}
\]

\end{parts}

\vspace{0.25em}

% ================================================================
% B ANSWERS
% ================================================================

\answerbar{B. Tangent Lines and Intersecting Curves}

\begin{parts}

\item
\[
\boxed{t=1,\ 3}
\]

\item
\[
\boxed{\cos\theta=\frac49}
\]

\end{parts}

\vspace{0.25em}

% ================================================================
% C ANSWERS
% ================================================================

\answerbar{C. Quadric Surfaces and Cross Sections}

\begin{parts}

\item
\[
\boxed{
\frac{(x-1)^2}{9}
+
\frac{(y+2)^2}{4}
-
\frac{z}{3}
=1
}
\]

\[
\boxed{\text{Elliptic paraboloid}}
\]

Axis of symmetry:
\[
\boxed{x=1,\qquad y=-2}
\]
which is parallel to the $z$-axis, and the surface opens in the
\[
\boxed{\text{positive }z\text{-direction}}.
\]

\item
\[
\boxed{
x=1:\text{ parabola},
\qquad
y=-2:\text{ parabola},
\qquad
z=0:\text{ ellipse}
}
\]

\item
\[
\boxed{6\pi}
\]

\end{parts}

\vspace{0.25em}

% ================================================================
% D ANSWERS
% ================================================================

\answerbar{D. Spherical and Cartesian Descriptions of a Surface}

\begin{parts}

\item
\[
\boxed{x^2+y^2+4z^2=8y}
\]

\item
\[
\boxed{
\frac{x^2}{16}
+
\frac{(y-4)^2}{16}
+
\frac{z^2}{4}
=1
}
\]

\[
\boxed{\text{Ellipsoid}}
\]

\item
Center:
\[
\boxed{(0,4,0)}
\]
with semiaxis lengths
\[
\boxed{4,\ 4,\ 2}.
\]

\end{parts}

\end{document}